\documentclass[11pt]{article}
% The preamble every project's paper shares: % The preamble every project's paper shares: % The preamble every project's paper shares: \input{../../preamble} from
% projects/<project>/paper/. Boards come from chessboard and xskak, plots from
% pgfplots; every number a paper prints comes from its generated.tex.
\usepackage[LSBC4,T1]{fontenc}
\usepackage[utf8]{inputenc}
\usepackage{lmodern}
\usepackage[margin=1in]{geometry}
\usepackage{microtype}
\usepackage{amsmath}
\usepackage{booktabs}
\usepackage{tabularx}
\usepackage{array}
\usepackage{xcolor}
\usepackage[obeyspaces,spaces]{url}
\usepackage{xskak}
\usepackage{pgfplots}
\pgfplotsset{compat=1.18}
\usepackage{caption}
\usepackage{subcaption}
\usepackage{enumitem}
\usepackage[section]{placeins}
\usepackage[hidelinks]{hyperref}

\setchessboard{boardfontencoding=LSBC4, boardfontfamily=skaknew, blackfieldcolor=black!22,
  whitefieldcolor=white, setfontcolors, showmover=false, label=true, labelfontsize=6pt,
  boardfontsize=14pt}

% A move or line in UCI, and a position in FEN, both breakable at a slash.
\newcommand{\uci}[1]{\texttt{#1}}
\newcommand{\fen}[1]{{\footnotesize\path{#1}}}
\newcolumntype{L}{>{\raggedright\arraybackslash}X}

\setlength{\parskip}{0.4em}
\setlength{\parindent}{0pt}
 from
% projects/<project>/paper/. Boards come from chessboard and xskak, plots from
% pgfplots; every number a paper prints comes from its generated.tex.
\usepackage[LSBC4,T1]{fontenc}
\usepackage[utf8]{inputenc}
\usepackage{lmodern}
\usepackage[margin=1in]{geometry}
\usepackage{microtype}
\usepackage{amsmath}
\usepackage{booktabs}
\usepackage{tabularx}
\usepackage{array}
\usepackage{xcolor}
\usepackage[obeyspaces,spaces]{url}
\usepackage{xskak}
\usepackage{pgfplots}
\pgfplotsset{compat=1.18}
\usepackage{caption}
\usepackage{subcaption}
\usepackage{enumitem}
\usepackage[section]{placeins}
\usepackage[hidelinks]{hyperref}

\setchessboard{boardfontencoding=LSBC4, boardfontfamily=skaknew, blackfieldcolor=black!22,
  whitefieldcolor=white, setfontcolors, showmover=false, label=true, labelfontsize=6pt,
  boardfontsize=14pt}

% A move or line in UCI, and a position in FEN, both breakable at a slash.
\newcommand{\uci}[1]{\texttt{#1}}
\newcommand{\fen}[1]{{\footnotesize\path{#1}}}
\newcolumntype{L}{>{\raggedright\arraybackslash}X}

\setlength{\parskip}{0.4em}
\setlength{\parindent}{0pt}
 from
% projects/<project>/paper/. Boards come from chessboard and xskak, plots from
% pgfplots; every number a paper prints comes from its generated.tex.
\usepackage[LSBC4,T1]{fontenc}
\usepackage[utf8]{inputenc}
\usepackage{lmodern}
\usepackage[margin=1in]{geometry}
\usepackage{microtype}
\usepackage{amsmath}
\usepackage{booktabs}
\usepackage{tabularx}
\usepackage{array}
\usepackage{xcolor}
\usepackage[obeyspaces,spaces]{url}
\usepackage{xskak}
\usepackage{pgfplots}
\pgfplotsset{compat=1.18}
\usepackage{caption}
\usepackage{subcaption}
\usepackage{enumitem}
\usepackage[section]{placeins}
\usepackage[hidelinks]{hyperref}

\setchessboard{boardfontencoding=LSBC4, boardfontfamily=skaknew, blackfieldcolor=black!22,
  whitefieldcolor=white, setfontcolors, showmover=false, label=true, labelfontsize=6pt,
  boardfontsize=14pt}

% A move or line in UCI, and a position in FEN, both breakable at a slash.
\newcommand{\uci}[1]{\texttt{#1}}
\newcommand{\fen}[1]{{\footnotesize\path{#1}}}
\newcolumntype{L}{>{\raggedright\arraybackslash}X}

\setlength{\parskip}{0.4em}
\setlength{\parindent}{0pt}

% Written by projects/findable-move/paper/generate.mjs; do not edit by hand.

\newcommand{\AppCommit}{f4bd63f19b}
\newcommand{\Offset}{250}
\newcommand{\Window}{3}
\newcommand{\Rollouts}{16}
\newcommand{\Plies}{12}
\newcommand{\HalfPlies}{six}
\newcommand{\HalfRollouts}{8}
\newcommand{\BudgetNodes}{120,000}
\newcommand{\Margin}{one}
\newcommand{\DecoyDrop}{10}
\newcommand{\NegMargin}{$-$1}
\newcommand{\Draws}{2,000}
\newcommand{\Games}{647}
\newcommand{\TrainGames}{431}
\newcommand{\HeldoutGames}{216}
\newcommand{\Examined}{2,013}
\newcommand{\WithTwo}{2,011}
\newcommand{\SameAsBest}{1,468}
\newcommand{\SameAsBestPct}{73}
\newcommand{\Entering}{543}
\newcommand{\Turning}{363}
\newcommand{\Plans}{180}
\newcommand{\TrainPositions}{352}
\newcommand{\HeldoutPositions}{191}
\newcommand{\RolloutCount}{22,784}
\newcommand{\DecoyTrainN}{189}
\newcommand{\DecoyHeldN}{113}
\newcommand{\Verdict}{unclear}
\newcommand{\DFTrain}{$-$0.93 ($-$2.18 to 0.28)}
\newcommand{\DFHeld}{$-$1.88 ($-$3.37 to $-$0.46)}
\newcommand{\RuleTable}{%
\begin{table}[ht]\centering\small
\begin{tabular}{lrr}
\toprule
 & Train & Held out \\
\midrule
Games & 228 & 135 \\
Positions & 352 & 191 \\
Null: B's two halves & $-$0.32 ($-$1.62 to 1.02) & 0.63 ($-$1.00 to 2.31) \\
Decoy: his errors against B, $D_P$ & $-$15.80 ($-$18.20 to $-$13.37) & $-$18.71 ($-$21.82 to $-$15.67) \\
\textbf{$D_F = V(F) - V(B)$} & \textbf{$-$0.93 ($-$2.18 to 0.28)} & \textbf{$-$1.88 ($-$3.37 to $-$0.46)} \\
\bottomrule
\end{tabular}
\caption{The controls and the verdict's measure, in win\% with 95\% intervals from 2,000 bootstrap draws over games. The decoy rows are the turning points where his move lost at least 10 win\% and is neither B nor F.}
\label{tab:rule}
\end{table}
}
\newcommand{\DFAll}{$-$1.27 ($-$2.21 to $-$0.34)}
\newcommand{\MinusDF}{1.27 (0.34 to 2.21)}
\newcommand{\RootGap}{1.01 (0.94 to 1.08)}
\newcommand{\FAheadPct}{44}
\newcommand{\PBestPct}{9}
\newcommand{\PFindPct}{33}
\newcommand{\FOverP}{14.64 (12.81 to 16.44)}
\newcommand{\FOverPN}{338}
\newcommand{\NullWidth}{2.1}
\newcommand{\SplitTable}{%
\begin{table}[ht]\centering\small
\begin{tabular}{lrrr}
\toprule
 & Positions & Games & $D_F$, win\% \\
\midrule
All & 543 & 363 & $-$1.27 ($-$2.21 to $-$0.34) \\
\midrule
B quiet & 448 & 321 & $-$0.85 ($-$1.85 to 0.16) \\
B forcing (a capture or a check) & 95 & 83 & $-$3.24 ($-$6.17 to $-$0.44) \\
\midrule
Turning points & 363 & 280 & $-$1.55 ($-$2.86 to $-$0.29) \\
Plan checkpoints & 180 & 164 & $-$0.68 ($-$2.05 to 0.66) \\
\midrule
Opening & 133 & 124 & $-$0.10 ($-$1.93 to 1.63) \\
Middlegame & 272 & 210 & $-$1.57 ($-$2.93 to $-$0.21) \\
Endgame & 138 & 104 & $-$1.78 ($-$3.92 to 0.11) \\
\bottomrule
\end{tabular}
\caption{$D_F$ over both sets, by the kind of the engine's best, by the kind of position and by the divider's phase. These splits were named before the run but carry no rule: they describe.}
\label{tab:splits}
\end{table}
}
\newcommand{\QuietN}{448}
\newcommand{\ForcingN}{95}
\newcommand{\HalvesMedian}{7.8}
\newcommand{\HalvesNinety}{21.0}
\newcommand{\QuietGap}{1.5}
\newcommand{\ForcingGap}{2}
\newcommand{\ExKind}{a plan checkpoint}
\newcommand{\ExPhase}{opening}
\newcommand{\ExSide}{Black}
\newcommand{\ExBest}{11\ldots{}\symking{}f8}
\newcommand{\ExBestUci}{\uci{e8f8}}
\newcommand{\ExFind}{11\ldots{}\symknight{}e7}
\newcommand{\ExFindUci}{\uci{g8e7}}
\newcommand{\ExPlayed}{11\ldots{}\symknight{}e7}
\newcommand{\ExPlayedUci}{\uci{g8e7}}
\newcommand{\ExPBest}{3}
\newcommand{\ExPFind}{79}
\newcommand{\ExGap}{2.0}
\newcommand{\ExVBest}{51.3}
\newcommand{\ExVFind}{76.1}
\newcommand{\ExDF}{24.9}
\newcommand{\ExFen}{\fen{r2qk1nr/1p3ppp/p1nb4/3p4/B7/5Q2/PPPP1PPP/R1B1R1K1 b kq - 1 11}}
\newcommand{\ExBoard}{%
\chessboard[setfen={r2qk1nr/1p3ppp/p1nb4/3p4/B7/5Q2/PPPP1PPP/R1B1R1K1 b kq - 1 11}, boardfontsize=14pt, pgfstyle=straightmove, arrow=to, linewidth=0.07em, color=blue!60!black, markmoves={e8-f8}, pgfstyle=straightmove, arrow=to, linewidth=0.07em, color=green!50!black, markmoves={g8-e7}]
}
\newcommand{\EyKind}{a turning point}
\newcommand{\EyPhase}{opening}
\newcommand{\EySide}{Black}
\newcommand{\EyBest}{8\ldots{}\symbishop{}xf3}
\newcommand{\EyBestUci}{\uci{g4f3}}
\newcommand{\EyFind}{8\ldots{}\symbishop{}b4}
\newcommand{\EyFindUci}{\uci{f8b4}}
\newcommand{\EyPlayed}{8\ldots{}\symknight{}xd4}
\newcommand{\EyPlayedUci}{\uci{c6d4}}
\newcommand{\EyPBest}{12}
\newcommand{\EyPFind}{29}
\newcommand{\EyGap}{1.9}
\newcommand{\EyVBest}{55.6}
\newcommand{\EyVFind}{42.0}
\newcommand{\EyDF}{$-$13.6}
\newcommand{\EyFen}{\fen{r2qkbnr/pp3ppp/2n1p3/3pP3/2PB2b1/2N2N2/PP3PPP/R2QKB1R b KQkq - 0 8}}
\newcommand{\EyBoard}{%
\chessboard[setfen={r2qkbnr/pp3ppp/2n1p3/3pP3/2PB2b1/2N2N2/PP3PPP/R2QKB1R b KQkq - 0 8}, boardfontsize=14pt, pgfstyle=straightmove, arrow=to, linewidth=0.07em, color=blue!60!black, markmoves={g4-f3}, pgfstyle=straightmove, arrow=to, linewidth=0.07em, color=green!50!black, markmoves={f8-b4}]
}


\title{Does a move he would find hold up as well as the engine's best?\\[0.3em]
\large Maia-2 plays on from the player's turning points and plan checkpoints}
\author{Chess Coach research lab}
\date{}

\begin{document}
\maketitle

\begin{abstract}
Game Review's better move is always the engine's top line. A human coach often shows another
move instead: one the player would find, nearly as good. In the player's 10-minute games, at
\SameAsBestPct\% of turning points and plan checkpoints the sound move players at his band would
most likely play is already the engine's best. In the other \Entering{} positions, each move was
played on \Rollouts{} times by Maia-2 at both players' bands for \Plies{} plies, and the ends were
scored by Stockfish. Both controls held. Against the engine's best, the move he would find scored
\DFTrain{} win\% on the train games and \DFHeld{} on the held-out games, so by the rule fixed
before the run the verdict is \emph{\Verdict}: a cost of more than \Margin{} win\% is neither
ruled out nor shown. Over both sets it cost \MinusDF{} win\% in human hands, against the
\RootGap{} the engine's own evaluation says it gives up: no hidden practical edge, and no hidden
trap. The cost is clearest when the engine's best is forcing and may be nothing when it is quiet.
At his turning points the move he would find scores \FOverP{} win\% above the move he played.
\end{abstract}

\section{The question}

Game Review names the engine's top line as the better move. A coach who knows the player
often shows a different one: a move he would find, close to the best, with a follow-up he can
play. The review's own window already treats a move within \Window{} win\% of the best as sound,
so such a move is never called an error. What the window cannot say is how the move fares when
people at his level play on from it. A move the engine rates a little lower may be easier to
handle, and score better in practice; or the engine's best may hold something the natural move
gives away, which shows only later.

In each position, three moves:

\begin{itemize}[nosep]
\item \textbf{B}, the engine's best;
\item \textbf{F}, the move he would find: the sound move Maia-2 rates most likely at his band;
\item \textbf{P}, the move he played.
\end{itemize}

The question is whether F, played on by people at these bands, scores as well as B.

\section{How it was measured}

\paragraph{The positions.} His 10-minute games (no increment) in the app's archive at commit
\texttt{\AppCommit}: \Games{} games, in order of their end time, the earlier two thirds
(\TrainGames) for training and the later third (\HeldoutGames) held out. Nothing is fitted, so
the two sets are two replications. Two kinds of position, both with him to move: the review's
turning points, with up to five engine lines from the routine's deep search, and its plan
checkpoints, with up to three. The candidates are the first moves of those lines, each with its
stored evaluation; a position needs two. A candidate is \emph{sound} when it gives up no more
than \Window{} win\% against B, the review's own window. His band is the Maia-2 level of his
rating at the game with the app's offset (+\Offset), and the opponent's band the same of the
opponent's rating. F is the sound candidate Maia-2 rates most likely at his band, conditioned on
the opponent's band.

Of the \Examined{} positions examined, \WithTwo{} have two candidates or more, and in
\SameAsBest{} of those (\SameAsBestPct\%) F is B: the move his band would most likely play among
the sound ones is the engine's best. The other \Entering{} enter the study: \Turning{} turning
points and \Plans{} plan checkpoints, \TrainPositions{} in the train games and
\HeldoutPositions{} held out.

\paragraph{The rollouts.} Each entering position plays out B and F, and P at a turning point when
it is neither. Each move gets \Rollouts{} rollouts. A rollout plays the move, then \Plies{} more
plies (\HalfPlies{} by each side), each sampled from Maia-2's policy for the side to move: his
band for his moves, the opponent's band for the opponent's, each conditioned on the other's.
Sampling is seeded by the position and the rollout. A rollout stops early when the game ends. Its
end is scored as his win\% by Stockfish through the app's engine service at the check budget
(\BudgetNodes{} nodes); a mate scores 100 or 0 and a draw 50. \RolloutCount{} rollouts in all.

$V(m)$ is the mean of a move's \Rollouts{} rollouts, and the measures are paired by position:
$D_F = V(F) - V(B)$, what the move he would find costs or gains in the hands of players at these
bands, and $D_P = V(P) - V(B)$, the same for his move. Intervals are the 95\% percentiles of
\Draws{} bootstrap draws over games.

\paragraph{The rule, fixed before any rollout was scored.} The controls come first, and if
either fails on either set there is no verdict. \emph{The null}: B's rollouts split into the
first \HalfRollouts{} and the last \HalfRollouts, and the mean difference's interval must cover
zero. \emph{The decoy}: at the turning points where his move lost at least \DecoyDrop{} win\%
and is neither B nor F, the interval of the mean $D_P$ must lie below zero; rollouts that cannot
see his known errors cannot judge F either. Then the verdict on F, from the mean $D_F$ on each
set: it \emph{holds up} when the interval's lower end is above \NegMargin{} win\% on both sets (\Margin{} win\%
is a third of the review's window); it is \emph{better in his hands} when the lower end is above
zero on both; \emph{the engine's best is worth the trouble} when the upper end is below \NegMargin{} on
both; and it is \emph{unclear} otherwise.

\section{The controls and the verdict}

\RuleTable

Both controls held on both sets (Table~\ref{tab:rule}). Two halves of the same move's rollouts
agree within noise, and the rollouts see his errors clearly: at his turning points where he lost
at least \DecoyDrop{} win\% (\DecoyTrainN{} train, \DecoyHeldN{} held out), his move scores far
below B in human hands too.

The move he would find scored a little below the engine's best on both sets. On the train games
the interval reaches above zero; on the held-out games it lies below zero but not below \NegMargin.
Neither set rules out a cost of more than \Margin{} win\%, and neither shows one, so the verdict
is \emph{\Verdict}.

\section{What the rollouts describe}

\SplitTable

These numbers describe; none of them carries a rule.

\paragraph{The cost is about what the engine says.} Over both sets the move he would find
gives up \MinusDF{} win\% in human hands. The engine's own evaluation of the same moves says it
gives up \RootGap. The interval of the first covers the second: in the hands of players at these
bands, the engine's gap neither closes nor opens. F scored above B in \FAheadPct\% of positions,
as a cost near zero with noisy rollouts would give.

\paragraph{F is the move people play.} At his band Maia-2 plays B \PBestPct\% of the time and F
\PFindPct\% of the time, on average over the entering positions.

\paragraph{The cost sits with forcing moves.} When B is quiet (\QuietN{} positions) the cost's
interval reaches zero. When B is a capture or a check (\ForcingN), F costs more, and the interval
lies below zero. The cost is larger at turning points than at plan checkpoints, and larger in the
middlegame and endgame than in the opening (Table~\ref{tab:splits}).

\paragraph{What costs him points is elsewhere.} At the \FOverPN{} turning points where he played
neither B nor F, F scores \FOverP{} win\% above the move he played. Choosing the natural sound
move over the engine's best costs about a point; missing both costs many.

\section{Two positions}

One position's rollouts are noisy: two halves of the same move's rollouts differ by a median of
\HalvesMedian{} win\%, and by \HalvesNinety{} or more in one position in ten. The two positions
below illustrate the summary; they do not test anything.

\begin{figure}[ht]
\centering
\begin{subfigure}[t]{0.48\textwidth}
\centering
\ExBoard\\[0.4em]
{\footnotesize FEN \ExFen}
\caption{\ExSide{} to move, \ExKind{} in the \ExPhase, in check. The engine's best (blue) is
\ExBest{} (\ExBestUci), which his band plays \ExPBest\% of the time; the move he would find
(green) blocks, \ExFind{} (\ExFindUci), at \ExPFind\%. The engine sees a gap of \ExGap{} win\%;
in the rollouts F scored \ExVFind{} and B \ExVBest, a difference of \ExDF. He played
\ExPlayed.}
\label{fig:held}
\end{subfigure}\hfill
\begin{subfigure}[t]{0.48\textwidth}
\centering
\EyBoard\\[0.4em]
{\footnotesize FEN \EyFen}
\caption{\EySide{} to move, \EyKind{} in the \EyPhase. The engine's best (blue) is the capture
\EyBest{} (\EyBestUci), which his band plays \EyPBest\% of the time; the move he would find
(green) is \EyFind{} (\EyFindUci), at \EyPFind\%. The engine sees a gap of \EyGap{} win\%; in
the rollouts B scored \EyVBest{} and F \EyVFind, a difference of \EyDF. He played \EyPlayed{}
(\EyPlayedUci).}
\label{fig:cost}
\end{subfigure}
\caption{Left: of the positions where both moves are quiet, the engine sees a gap of at least
\QuietGap{} win\% and F scored at least as well as B, the one where Maia-2 rates F most above B.
Right: of the positions where B is forcing, F is quiet and the engine's gap is at most
\ForcingGap{} win\%, the one where F scored furthest below B.}
\label{fig:examples}
\end{figure}

\section{What the app can take}

\begin{itemize}[nosep]
\item The review can show the move he would find beside the engine's best when the best is quiet,
  as what it is: a sound move he would find, costing about what the review's window says. The
  study does not show it is better.
\item When the engine's best is a capture or a check, the review should keep pointing to it:
  there the move he would find costs more in human hands.
\item The larger gain is in the moves he plays that are neither. Coaching him toward a sound move
  he would find, whichever it is, is worth far more than coaching him from that move to the
  engine's.
\end{itemize}

\section{Caveats}

\begin{itemize}[nosep]
\item \textbf{Maia-2 plays both sides.} The rollouts model players at these bands, not him or
  this opponent. A model of him is a later study.
\item \textbf{\Plies{} plies.} The end is scored by Stockfish, so what comes after counts only as
  far as the engine sees it then.
\item \textbf{Only the stored lines.} F is the most natural of the two to five moves the review
  kept, not of every legal move.
\item \textbf{Near a decided game the window is wide.} Win\% flattens close to a win or a loss, so
  moves the window admits can differ by more than their win\% shows; the rollouts see it, the
  window does not.
\item \textbf{The splits describe.} A second round could test quiet best moves alone, on positions
  this round did not see, with its margin fixed first.
\end{itemize}

\end{document}
